% !PRINTED 203
\weberpage{203}
%<<CONT
bestimmen kann, so kann daraus $\alpha^2$ etwa aus
\[
\frac{\alpha\beta\,.\,\alpha\gamma}{\beta\gamma}
\]
berechnet werden.

Es kommt also jetzt noch darauf an, die cubische Gleichung
zu bilden, deren Wurzeln $a_2', a_2'', a_2'''$ sind. Diese erhält man
aber sofort aus den Invarianten.

Wenn man die Invarianten $A', B'$ aus der transformirten
Form (4) bildet, so ergiebt sich nach §.\,64, (8), (9), (12), da
$r = 1$ ist,
\weq{11}{\begin{gathered}
A = a_2'^{\,2} - 3\,a_1'\,a_3',\\
B = 2\,a_2'^{\,3} - 9\,a_1'\,a_2'\,a_3',\\
D = a_1'^{\,2}\,a_2'^{\,2}\,a_3'^{\,2} - 4\,a_1'^{\,3}\,a_3'^{\,3},
\end{gathered}}
und daraus durch Elimination von $a_1'\,a_3'$
\[
a_2'^{\,3} - 3\,a_2'\,A + B = 0.
\]
Daher sind $a_2', a_2'', a_2'''$ die Wurzeln der für $z$ cubischen
Gleichung
\weq{12}{z^3 - 3\,A\,z + B = 0.}

Dies ist wohl die einfachste Form, die man der cubischen
Resolvente der biquadratischen Gleichung geben kann.

\begin{center}
§. 66.

\gesperrt{Die Covarianten.}
\end{center}

Wenn wir die Covariante $H$, nach §.\,64, (2) für die transformirte
 Form
\weq{1}{f(x,\,y) = F(\xi,\,\eta) = a_1'\,\xi^3\eta + a_2'\,\xi^2\eta^2 + a_3'\,\xi\eta^3}
bilden, so erhalten wir
\weq{2}{\begin{gathered}
-H = 3\,a_1'^{\,2}\,\xi^4 + 4\,a_1'\,a_2'\,\xi^3\eta - (6\,a_1'\,a_3' - 4\,a_2'^{\,2})\,\xi^2\eta^2\\
{} + 4\,a_2'\,a_3'\,\xi\eta^3 + 3\,a_3'^{\,2}\,\eta^4,
\end{gathered}}
und daraus folgt
\weq{3}{H + 4\,a_2'\,f = -3\,(a_1'\,\xi^2 - a_3'\,\eta^2)^2.}

Es ist also diese Verbindung, von dem Factor $-3$ abgesehen,
das Quadrat einer quadratischen Form, und dasselbe gilt, wenn
wir $a_2'$ durch $a_2'', a_2'''$ ersetzen. Wir setzen also zur
Abkürzung
%CONT>>
